\documentclass{report}  
\usepackage[utopia]{mathdesign} 
%\usepackage{amsmath,amsfonts,amsthm,amssymb,mathtools}

\input{preamble.tex}

\usepackage{amsmath,amsthm,mathtools}
\usepackage{titlesec}
\usepackage{microtype}
\definecolor{myblue}{RGB}{0,82,155}




%\usepackage[scr]{rsfso}

%\usepackage{libertine}
%\usepackage{mathpazo}
%\usepackage{palatino}
%usepackage{crimson}


\newcommand{\faketarget}{\oplus\!\!\!\!\odot}
\newcommand{\target}{%
  \begin{tikzpicture}[scale=0.5]
    \fill[black] (0,0) circle (0.1);
    \draw (0,0) circle (0.2);
    \draw (0,0) circle (0.3);
  \end{tikzpicture}%
}

\usepackage[clock]{ifsym}


\title{\Huge{Structure de données}\\{IFT2015}\\{\textbf{Introduction}}}
\author{\huge{Franz Girardin}}
\date{\today}
\lstset{inputencoding=utf8/latin1}

            %%%%%%%%%%%%%%%%%  Sect.                          %%%%%%%%%%%%%%%%%%%%%%%%%%%%%%%%%%%%%%%%%%%%%%%%%%%%%%%%%
\usepackage{helvet}
\titleformat{\chapter}[display]
  {\normalfont\bfseries\color{myblue}}
  {\filleft%
    \begin{tikzpicture}
    \node[
      outer sep=0pt,
      text width=1.5cm,
      minimum height=2cm,
      fill=myblue,
      font=\color{white}\fontsize{40}{50}\selectfont,
      align=center
      ] (num) {\thechapter};
    \node[
      rotate=90,
      anchor=south,
      font=\color{black}\Large\normalfont
      ] at ([xshift=-5pt]num.west) {\textls[180]{\textsc{Section}}};  
    \end{tikzpicture}%
  }
  {5pt}
  {\titlerule[2.0pt]\vskip3pt\titlerule\vskip4pt\large\normalfont}

\titleformat{\section}
  {\normalfont\scshape}{\thesection}{1em}{}



\titleformat{\section}
  {\normalfont\scshape}{\thesection}{1em}{}


% Customizing the spacing for the chapter titles
\titlespacing*{\chapter}{0pt}{0pt}{20pt}

%\usepackage[utopia]{mathdesign}
% Allow hfill in math environment
\newcommand{\specialcell}[1]{\ifmeasuring@#1\else\omit$\displaystyle#1$\ignorespaces\fi}

% Allow you to do the non implication (implication barred)
\newcommand{\notimplies}{%
  \mathrel{{\ooalign{\hidewidth$\not\phantom{=}$\hidewidth\cr$\implies$}}}}



\DeclareRobustCommand{\looongrightarrow}{%
  \DOTSB\relbar\joinrel\relbar\joinrel\relbar\joinrel\rightarrow
}


\title{\Huge{Interface PM}\\{IFT2905}\\{\textbf{Fuille de notes}}}
\author{\huge{Franz Girardin}}
\date{\today}


\begin{document}
\maketitle
\pagebreak
\tableofcontents
\pagebreak
\begin{multicols*}{3}


    \footnotesize
    %\section*{Design of everyday things}

    \chapter{Introduction aux structures de données}
    \section{Caractéristique génériques de $\mathbb{SD}$ }

    \paragraph{Comment décrire une structure de données}
    \begin{itemize}
      \item [$\rhd$ ] \textbf{Linéaire ou \textit{non} linéaire  }  
        \begin{itemize}
          \item [$\blacktriangleright$ ] P. ex. 
            \texttt{graph} vs. \texttt{arrays}    
        \end{itemize}
      \item [$\rhd$ ] Homogène ou non homogène  
        \begin{itemize}
          \item [$\blacktriangleright$ ] Indique si les données sont mê 
            \texttt{type}  
        \end{itemize}
      \item [$\rhd$ ] Statique ou \textit{dynamique}  
        \begin{itemize}
          \item [$\blacktriangleright$ ] P. ex. taille modifiable ou 
            \texttt{fixe}.   
        \end{itemize} 
    \end{itemize}
    
    \begin{note}{}{}
        Plusieurs problèmes ont une complexité telle que les $\mathbb{SD}$ 
        \textbf{base} sont inappropriées, d’où la nécessité 
        des $\mathbb{SD}$ \textbf{abstraites}, plus adéquates.  
    \end{note}

    \paragraph{$\mathbb{SD}$ Abstaites}
    Elles \texttt{généralise} les types de données de base et permettent 
    l'émergence de concept sophistiqués d'\textbf{abstraction}, 
    d'\textbf{encapsulation} et de typage. Elle possèdent 
    \textbf{5 caractéristiques} fondamentales ($\mathbb{TUOPA}$): 
    \begin{itemize}
      \item [$\rhd$ ]  $\mathbb{T}$ype
      \item [$\rhd$ ] $\mathbb{U}$tilise 
      \item [$\rhd$ ] $\mathbb{O}$pérations
      \item [$\rhd$ ] $\mathbb{P}$rocédures
      \item [$\rhd$ ] $\mathbb{A}$xiomes 
    \end{itemize}
    
    \section{Tableau}
    \noindent
    Stocke $\texttt{E}$ dans emplacement de mémoire \textcolor{red}{contigus}.
    \begin{itemize}
      \item [$\rhd$ ] De longueur fixe ou variable
        \begin{itemize}
          \item [$\blacktriangleright$ ] Autrefois : adresse brutes
          \item [$\blacktriangleright$ ] Maintenant : $N$ \texttt{arrays} (indexation)  
          \item [$\blacktriangleright$ ] Maintenant : \texttt{records} (champs)  
        \end{itemize}
    \end{itemize}

    \section{Pile}
    La propriété fondamentale d'une pile est que le seul éléement accessible 
    et celui se trouvant au sommet de la pile. La \textbf{pile}  suit le principe 
    \textbf{LIFO} : \textit{last in, first out}.     



    \section{File}
    La \textbf{file}   suit le principe 
    \textbf{FIFO} : \textit{first in, first out}. Elle possède, entre autres,
    les opération enfiler \textit{enqueue} et défiler \textit{dequeue}.      

    \section{Liste}
    Une file est le résultat de l'\textit{ordonancement} d'un nombre 
    \textcolor{red}{dénombrable} de données où la même valeur 
    peut apparaître plusieurs fois. 

    \begin{itemize}
      \item [$\rhd$ ] \textbf{Implémentation}  
        \begin{itemize}
          \item [$\blacktriangleright$ ] Liste chaînée \texttt{Linked List}  
          \item [$\blacktriangleright$ ] Liste doubleemnt chaînée \texttt{Linked List}  
          \item [$\blacktriangleright$ ] Liste circulaire 

          \item [$\blacktriangleright$ ] Tableau \texttt{array}  
        \end{itemize}
    \end{itemize}


    \begin{note}{}{}
        Une liste chaînée stocke un ensemble d’éléments de 
        façon linéaire. Chaque élément ou nœud d’une liste chaînée 
        contient un \textbf{élément de données} ainsi qu’une 
        \textbf{référence}, ou lien, 
        vers l’élément suivant de la liste.       
    \end{note}


    \section{Arbre} 
    Un arbre stocke un ensemble d’éléments sous une
    forme \textbf{hiérarchique abstraite}. Chaque nœud est relié
    aux autres et peut contenir plusieurs sous-valeurs
    appelées enfants. Il s'agit d'un \texttt{graphe} avec 
    \textbf{3 particularité fondamentales}: 
    \begin{itemize}
      \item [$\rhd$ ] Acyclique
      \item [$\rhd$ ] Connexe
      \item [$\rhd$ ] Possède une seule racine
    \end{itemize}


    \section{File de priorité}
    Possède les mêmes propriété de base qu'un file, sauf que chaque 
    élément a en plus un poids de priorité qui détermine l'ordre 
    global des éléments. 

    \section{Graphe}
    Un graphe stocke un ensemble d’éléments de façon
    non linéaire. Il se compose d’un ensemble fini de
    nœuds, appelés \textbf{sommets}, et de lignes, les arêtes, qui
    relient les sommets entre eux. Les graphes permettent 
    notamment de représenter des systèmes réels, 
    comme des réseaux informatiques.

    \paragraph{Théorème des quatre couleurs}
    Il existe une $4$-coloration de n'importe qu'elle pays 
    qui a été découpé sur une carte. 


    \chapter{Pile}
    \paragraph{Exemples d'applications}
    \begin{itemize}
      \item [$\rhd$ ] Historique de pages visitées
      \item [$\rhd$ ] Annuler une séquence de texte dans un éditeur 
      \item [$\rhd$ ] Chaîne de méthodes appelées d'un fonction complexe
    \end{itemize}

    \paragraph{5 Opération fondamentales}
    \begin{itemize}
      \item [$\rhd$ ] Créer un pile 
      \item [$\rhd$ ] Empiler 
      \item [$\rhd$ ] Dépiler 
      \item [$\rhd$ ] Regarder le sommet 
      \item [$\rhd$ ] Obtenir le nombre d'éléments
    \end{itemize}

    \begin{lstlisting}
elements[0...n-1] <-- tableau // tableau taille n
// top <-- 0 pour assigner 0 à top on fait :

push(E) :
element[top] <-- E
top <-- top + 1 // modifie l'index du sommet 


//  Pour obtenir l'E au sommet, on fait : 

pop() :
retourne element[top] 
top <-- top -1 //decremente la valeur du sommet
x <-- element[top + 1] // enregistre val sommet  
element[top + 1] <-- null
    \end{lstlisting}


    \paragraph{Conventions d'écriture}

    \begin{itemize}
      \item [$\rhd$ ] \textbf{IN}fixé : \texttt{gauche \textcolor{red}{racine} droite}   
        \begin{itemize}
          \item [$\blacktriangleright$ ] 
            $5 \; \textcolor{red}{ - } \;   6 \textcolor{red}{\times} 7
            \leftrightarrow (5 - 6) \times 7 $  
        \end{itemize}
      \item [$\rhd$ ] \textbf{POST}fixé : 
        \texttt{gauche droite \textcolor{red}{racine}}   
        \begin{itemize}
          \item [$\blacktriangleright$ ] 
            $56 \textcolor{red}{ - }   7 \textcolor{red}{\times} 
            \leftrightarrow (56-) \; 7\times$  
        \end{itemize}
      \item [$\rhd$ ] \textbf{PRÉ}fixé : \texttt{\textcolor{red}{racine} gauche droite}  
        \begin{itemize}
          \item [$\blacktriangleright$ ] 
            $\times \textcolor{red}{ - }567 \textcolor{red}{\times} 
            \leftrightarrow \times (-56) \; 7$  
        \end{itemize}
    \end{itemize}













    \end{multicols*}


\end{document}

